\exnumber{01}
\extitle{UI III - Ex 2 - États multiples}
\exfiles{ex01/fichier Figma, nomme Name\_FirstName\_DeliverableName\_Date\_VersionNumber}
\exauthorize{Those named in the exercise}
\makeheaderfiles

\textbf{Instructions}:

Go to Sony.com \url{http://sony.com} and observe the buttons and call to actions. When you hover or click on them, they change!

Make screenshots of the different states of buttons, menu titles, filters etc. Then, gather and organize them.

Find at least five combinations of multi state elements.

\textbf{Livrables :}

\begin{itemize}
\item A Figma file with 5 combinations of multi state elements.
\item Déposez le fichier Figma sur Github, nommé ainsi : « Name\_FirstName\_DeliverableName\_Date\_VersionNumber ». Exemple : « Doe\_John\_MidFiPrototype\_05242024\_V1 ».
\item Ajoutez un titre dans le fichier.
\item Soignez l'aspect global du fichier. Il doit être clair, simple et lisible. Vous pouvez regarder canva.com, sans surcharger le dessin de détails.
\end{itemize}

\textbf{Ressources :}

\begin{itemize}
\item Video tutorial on Figma \url{https://www.youtube.com/watch?v=FTFaQWZBqQ8}
\item Beginning graphic design \url{https://edu.gcfglobal.org/en/beginning-graphic-design/}
\item Gareth graphic studio \url{https://www.youtube.com/c/GarethDavidStudio}
\item Example here with Nintendo.com \url{https://www.figma.com/file/ahSnWRZeKBO62oJDiXltxY/UI-III---Ex-2}
\end{itemize}

/!\textbackslash{} The screenshots date back to nov 2021. They may differ from the current version of Nintendo.com \url{http://Nintendo.com} but the rules are the same

